% Transcription — fonds Alexandre Grothendieck, shelfmark 158, batch 5 (pages 81-85).
% Established from the facsimile archives/batches/158/batch-05.pdf (a mirror of
% https://grothendieck.umontpellier.fr/158.pdf), per the skill
% transcribe-grothendieck. The batch file carries no cover sheet: PDF page 1 =
% archivists' page 81, checked against the pencilled numbers bottom-left, which
% read 81 on PDF page 1 and 85 on PDF page 5 — offset zero across the whole batch.
% Pages 82, 84 and 85 carry no mathematics — absent.
% Reconciled against the folder's other batches, 2026-09-21: the argument of E
% and R, first read as n, is restored to x on the evidence of leaf 74 and of
% this leaf's own E(x)'. The two underbrace sources are left as drawn.
% Pass: Opus 5 (claude-opus-5), 2026-09-20 — first pass, unchecked against the pages by a human.
% Opus 5 rather than the skill's Fable 5.1 default, for continuity with folders
% 159 and 160 of the same inventory group, both transcribed on Opus 5.
\documentclass[11pt,a4paper]{article}
% Le préambule commun à toutes les transcriptions du fonds Grothendieck.
%
% Il tient en une idée : **l'incertitude se compose, elle ne se résout pas.**
% Une transcription qui lisse un mot douteux a détruit la seule chose pour
% laquelle on la faisait — un lecteur doit pouvoir savoir, à chaque endroit,
% ce qui était sur le feuillet et ce qui a été ajouté. Les six macros qui
% suivent sont donc l'appareil critique, et non de la décoration.
%
% Les mêmes commandes sont reconnues par `scripts/render.mjs`, qui produit la
% vue de lecture du site. Une macro ajoutée ici doit l'être aussi là-bas,
% faute de quoi le rendu échouera bruyamment — ce qui est voulu : un rendu
% silencieusement faux est pire qu'un rendu absent.

\ProvidesPackage{grothendieck}[2026/08/08 Transcriptions du fonds Grothendieck]

\usepackage{amsmath,amssymb,amsthm}
\usepackage{mathrsfs}
\usepackage{stmaryrd}
% Les diagrammes commutatifs sont partout dans ce fonds : c'est la langue
% même de la géométrie algébrique de Grothendieck, et une transcription qui
% les rendrait en prose ne transcrirait rien.
\usepackage{tikz-cd}
% Français par défaut — c'est la langue des feuillets — mais l'anglais est
% chargé aussi : la traduction et le résumé sont des documents anglais, et la
% typographie française (espaces avant les deux-points) y serait une faute.
% Ces documents basculent par \selectlanguage{english} en tête de corps.
\usepackage[english,french]{babel}
\usepackage{fontspec}
\usepackage{geometry}
\geometry{a4paper,margin=25mm,marginparwidth=28mm}
\usepackage{marginnote}
\usepackage{xcolor}
% ulem, not soul. `soul`'s \st and \ul are built on a character-by-character
% scan that XeLaTeX's font machinery breaks: the document fails with an
% undefined control sequence rather than degrading. ulem does the same two jobs
% and is Unicode-engine safe. `normalem` stops it hijacking \emph, which the
% transcriptions use for Grothendieck's own emphasis.
\usepackage[normalem]{ulem}
\usepackage{hyperref}
\hypersetup{colorlinks=true,linkcolor=black,urlcolor=black,pdfborder={0 0 0}}

% --- Métadonnées du lot -----------------------------------------------------
% Reprises telles quelles par le script de rendu, qui les affiche en tête de
% la vue de lecture. Elles fixent ce que le fichier prétend couvrir : sans
% elles, une transcription flotte sans point d'attache dans un fonds de seize
% mille feuillets.

\newcommand{\folder}[1]{\gdef\grothFolder{#1}}
\newcommand{\batch}[1]{\gdef\grothBatch{#1}}
\newcommand{\leaves}[2]{\gdef\grothFirst{#1}\gdef\grothLast{#2}}
\newcommand{\foldertitle}[1]{\gdef\grothFoldertitle{#1}}
\gdef\grothFolder{?}\gdef\grothBatch{?}\gdef\grothFirst{?}\gdef\grothLast{?}\gdef\grothFoldertitle{}

% --- Le résumé --------------------------------------------------------------
%
% Il ouvre la lecture modernisée, sous le titre « Résumé », et il est composé
% en petit. Ce n'est pas un ornement : c'est le seul passage du document qui
% ne soit pas des mathématiques, et un lecteur doit pouvoir le reconnaître
% avant de l'avoir lu — puis le sauter s'il connaît déjà le sujet, ou s'y
% arrêter s'il ne le connaît pas. Le filet qui le suit marque la reprise du
% corps.

\newenvironment{resume}
  {\small\setlength{\parskip}{0.45em}}
  {\par\medskip\hrule\medskip}

% --- Mots-clés ---------------------------------------------------------------
%
% En anglais, en clôture du résumé de la lecture modernisée : le vocabulaire
% moderne sous lequel chercher ce que les feuillets construisent. Le manifeste
% du site (`npm run manifest`) les extrait pour étiqueter la cote entière —
% cette ligne est la source unique des tags, il n'y a pas de fichier à côté.
\newcommand{\keywords}[1]{\par\smallskip\noindent
  {\footnotesize\textsc{Keywords} --- \textit{#1}}\par}

% --- L'appareil critique ----------------------------------------------------

% Le numéro de feuillet, en marge. C'est l'ancre qui permet de régler un
% désaccord sur une formule en regardant *un* feuillet plutôt qu'une chemise
% de six cents. Le site s'en sert aussi pour tourner les pages du fac-similé
% quand on fait défiler la transcription.
\newcommand{\leaf}[1]{%
  \marginnote{\footnotesize\color{gray!70}\textsf{#1}}%
  \label{leaf:#1}\ignorespaces}

% Illisible : on ne devine pas. Un mot inventé qui se lit comme les autres est
% le pire résultat possible.
\newcommand{\ill}{\textcolor{red!60!black}{[\,\ldots\,]}}

% Lecture douteuse : le mot est donné, et signalé comme incertain.
\newcommand{\uncertain}[1]{\uline{#1}}

% Ajout de l'éditeur — une lettre manquante, un mot sous-entendu.
\newcommand{\add}[1]{\textcolor{blue!50!black}{[#1]}}

% Biffé par Grothendieck lui-même. Ce qu'il a rayé fait partie du document :
% une rature dit souvent où la pensée a changé de direction.
\newcommand{\struck}[1]{\sout{#1}}

% Note du transcripteur, sur l'état matériel du feuillet.
\newcommand{\note}[1]{\par\smallskip{\small\color{gray!60!black}\textit{[#1]}}\par\smallskip}

% Note marginale de Grothendieck, distincte du corps du texte.
\newcommand{\marginal}[1]{\par\smallskip\noindent\hspace{1em}%
  {\small\color{gray!60!black}#1}\par\smallskip}

% --- Titre ------------------------------------------------------------------

\AtBeginDocument{%
  \begin{center}
    {\large\bfseries Fonds Alexandre Grothendieck --- cote n\textsuperscript{o}~\grothFolder}\\[2pt]
    {\itshape\grothFoldertitle}\\[4pt]
    {\small feuillets \grothFirst--\grothLast\ (lot \grothBatch)}\\[2pt]
    {\footnotesize\color{gray!60!black}Transcription établie d'après le fac-similé numérisé
     par l'Université de Montpellier.\\
     Les lectures incertaines sont soulignées, les illisibles notées [\,\ldots\,],
     les ajouts entre crochets.}
  \end{center}
  \vspace{1em}}

\endinput


\folder{158}
\batch{5}
\pages{81}{85}
\dating{[vers 1990-1991]}
\watermark{Édition de démonstration}
\foldertitle{[Autour des Dérivateurs] : notes manuscrites (s.d.)}

\begin{document}

\section*{[Les quatre composés $\lambda v$, $\mu w$, $\lambda' v$, $\mu' w$ et la
fibre $\widetilde{X}_{/\!/(\tilde a,\tilde b)}$]}

\note{le titre ci-dessus est de l'éditeur : le lot est la queue du dossier, ses
deux feuillets écrits ne portent pas de titre et reprennent au milieu d'un
raisonnement commencé dans les lots précédents. Les deux feuillets sont paginés
12 et 13 de sa main et se suivent immédiatement}

\note{notations tenues pour tout le lot : le tilde marque le relevé d'un objet de
$X$ dans $\widetilde{X}$ ($\tilde x$, $\tilde a$, $\tilde b$, $\tilde c$) ;
$X_{/\!/(a,b)}$ est écrit avec une double barre oblique ; $R$ est
la relation par laquelle il divise les Hom, et son argument est une seule lettre,
lue $n$ ; $\sim$ est la congruence mod $R$, $=$ l'égalité stricte, et l'opposition
des deux est tout le sujet du lot}

\page{81}
\note{feuillet paginé 12 de sa main, en haut au centre. Il s'ouvre sur un premier
tracé du diagramme qui suit, entièrement biffé}

Il faut interpréter $\widetilde{X}_{/\!/(\tilde a,\tilde b)}$, \uncertain{disons}
dans

\begin{tikzcd}[column sep=small, row sep=small]
 & \tilde a \arrow[dr, dashed] & \\
\tilde x \arrow[ur, "v"] \arrow[dr, "w"'] & & \tilde c \\
 & \tilde b \arrow[ur, dashed] &
\end{tikzcd}

\note{le tracé est très rapide : les deux flèches issues de $\tilde a$ et de
$\tilde b$ sont en pointillé et leur but, lu $\tilde c$, reste douteux}

\ill{} au-dessus des

\begin{tikzcd}[column sep=small, row sep=small]
 & a \\
x \arrow[ur, "v"] \arrow[dr, "w"'] & \\
 & b
\end{tikzcd}

\note{quatre flèches partent encore de $a$ et de $b$ vers la droite, étiquetées
$\lambda$, $\lambda'$, $\mu$, $\mu'$ ; leurs buts, tracés dans un fouillis de
traits, ne se lisent pas et ne sont pas portés sur le diagramme}

Donc $\tilde x$ \struck{\ill{}} $E(x)/R(x)$ peut être défini en définissant par
\note{l'argument de $E$ et de $R$ est ici un seul jambage, lu d'abord $n$ ;
on le rétablit en $x$, d'une part parce que le même feuillet écrit plus bas
$E(x)'$ du même tracé, d'autre part parce que le feuillet 74 porte la même
construction en toutes lettres : $E(x) = \mathrm{Hom}(x,\xi)\amalg\mathrm{Hom}(x,\xi')$
modulo $R_{(x)}$, engendrée par les deux mêmes relations}
\[
  \underbrace{\lambda v,\; \mu w}_{\mathrm{Hom}(m,\xi)}
  \qquad
  \underbrace{\lambda' v,\; \mu' w}_{\mathrm{Hom}(a,\xi')}
\]
\note{les deux arguments sous les accolades sont lus $m$ et $a$ ; le premier est
tracé comme un $m$, le second comme un $a$, alors que la marche du raisonnement
demanderait la même source de part et d'autre. La lecture est laissée telle
quelle : le feuillet 74 écrit la même somme avec une seule source,
$\mathrm{Hom}(x,\xi)\amalg\mathrm{Hom}(x,\xi')$, mais le tracé d'ici ne la
porte pas et on ne le corrige pas sur la foi d'un autre feuillet}

dans $E(x)'$ (disons leurs images dans $\mathrm{Hom}/R(x)$) i.e. ces quatre il
\ill{} \uncertain{toutes} congrues mod $R(x)$. On le sait déjà pour
$\lambda v \sim \lambda' v$ et $\mu w \sim \mu' w$, en fait
$\tilde x = v^{*}(\tilde a) = w^{*}(\tilde b)$, donc il faut comparer l'égalité
\[
  v^{*}(\tilde a) = w^{*}(\tilde b)
\]
\struck{\ill{}} \uncertain{en cours}
\[
  \boxed{\;\lambda v \sim \mu w\;}
\]
\note{la dernière formule est encadrée sur le feuillet ; le feuillet s'arrête là}

\page{83}
\note{feuillet paginé 13 de sa main ; il fait suite immédiate au précédent}

Malheureusement, on ne peut en conclure $\lambda v = \mu w$, p. ex. si on a
$\lambda' v = \mu' w$, on \uncertain{en conclura} $\lambda v \sim \mu w$, mais
nullement $\lambda v = \mu w$. Plus précisément, dès qu'on se donne dans $X$ des
flèches $v$, $w$ \struck{\ill{}}

\begin{tikzcd}[column sep=small, row sep=small]
 & a \\
x' \arrow[ur, "v"] \arrow[dr, "w"'] & \\
 & b
\end{tikzcd}

avec $\lambda' v = \mu' w$, on en \uncertain{conclut} \ill{} des $\widetilde{X}$
un $\tilde x$ \ill{} $(\tilde a, \tilde c, \tilde b)$ \ill{}, mais rien \ill{}
\struck{\ill{}} $\lambda v = \mu w$.

ainsi dans le \ill{} de. On \ill{} \struck{\ill{}} \uncertain{pour} $f$,
\note{une addition interlinéaire de sa main, rattachée par un long trait courbe à
cet endroit de la ligne : « \uncertain{savait déjà qu'il} »}
--- (par \uncertain{tr.} \ill{}) $v$, $w$ de façon : rendre \ill{}
\note{le mot est souligné par lui}
des couples pertinents \ill{} (disons $\lambda' v = \mu' w$), de que nous
\uncertain{savons} déjà
\[
  \widetilde{X}_{/\!/(\tilde a,\tilde b)} \neq \emptyset
\]
--- mais je ne vois pas de moyen \ill{} montrer en $=$ temps que l'autre le
$\tilde x$. ---
\note{« en $=$ temps » : il écrit le signe $=$ pour \emph{même}. Le feuillet, et
avec lui le dossier, s'arrête sur cette phrase}

\end{document}
